\documentclass[12pt, a4paper]{article}

\usepackage[utf8]{inputenc}
\usepackage[T1]{fontenc}
\usepackage[french]{babel}
\usepackage{amsmath}
\usepackage{amssymb}

\usepackage{geometry}
\geometry{margin=2.5cm}
\usepackage{parskip}
\usepackage{microtype}

\usepackage{hyperref}
\hypersetup{
    colorlinks=true,
    linkcolor=black,
    urlcolor=blue,
    citecolor=black
}

% --- En-têtes ---
\usepackage{fancyhdr}
\pagestyle{fancy}
\fancyhf{}
\rhead{Conception de la base de données --- Ciqly}
\lhead{}
\cfoot{\thepage}

\usepackage{mdframed}

\newmdenv[
    linewidth=0.5pt,
    linecolor=black,
    leftmargin=0pt,
    rightmargin=0pt,
    innerleftmargin=10pt,
    innerrightmargin=10pt,
    innertopmargin=8pt,
    innerbottommargin=8pt
]{rappelbox}


\newmdenv[
    linewidth=0.4pt,
    linecolor=black,
    leftmargin=8pt,
    rightmargin=0pt,
    innerleftmargin=8pt,
    innerrightmargin=8pt,
    innertopmargin=6pt,
    innerbottommargin=6pt,
    topline=false,
    bottomline=false,
    rightline=false
]{notebox}

\newcommand{\attr}[1]{\texttt{#1}}
\newcommand{\citeg}[1]{\textnormal{\footnotesize[#1]}}

% ------------------------------------------------------------------
\begin{document}

    \begin{titlepage}
        \centering
        \vspace*{3cm}
        {\LARGE \textbf{Conception de la base de données}}\\[1em]
        {\large Base de donnée du site web Ciqly}\\[3cm]
        \vfill
    \end{titlepage}

    \tableofcontents
    \newpage

% ==================================================================


    \section{Rappels}

    \begin{rappelbox}
        \textbf{Notion XVII.5 : Dépendance fonctionnelle (\textit{Functional dependency})}
        \hfill\citeg{Gardarin, p.\,685}\\[0.5em]
        Soit $R(A_1, A_2, \ldots, A_n)$ un schéma de relation, et $X$ et $Y$ des sous-ensembles de
        $\{A_1, A_2, \ldots, A_n\}$. On dit que $X \rightarrow Y$ ($X$ détermine $Y$, ou $Y$ dépend
        fonctionnellement de $X$) si, pour toute extension $r$ de $R$, pour tout tuple $t_1$ et $t_2$ de
        $r$, on a :
        \[
            \Pi_X(t_1) = \Pi_X(t_2) \;\Rightarrow\; \Pi_Y(t_1) = \Pi_Y(t_2)
        \]
    \end{rappelbox}

    \begin{notebox}
        \small
        \textbf{Remarque.} Une dépendance fonctionnelle est une \emph{assertion sur toutes les valeurs
        possibles} des attributs, et non sur leurs valeurs actuelles : elle caractérise une intention et non
        une extension de la relation. Il est donc impossible de déduire les DF d'une réalisation particulière ;
        leur identification exige une analyse sémantique du domaine métier.
        \hfill\citeg{Gardarin, p.\,685}
    \end{notebox}

% ==================================================================


    \section{Dépendances fonctionnelles}

    On note respectivement $U$ et $V$ les relations comprenant l'ensemble des attributs nécessaires à
    garantir les informations fournies par les deux fichiers de données de l'ANSES : la table CIQUAL
    et la table des aliments moyens. Une étude approfondie a été menée tant sur les données \texttt{XML}
    et \texttt{XLSX} que sur les fichiers de documentation mis à disposition.

% ------------------------------------------------------------------

    \subsection{Dépendances fonctionnelles de \texorpdfstring{$U$}{U}}

    À partir de la documentation fournie et des fichiers XML, on déduit les dépendances fonctionnelles
    afin de constituer un schéma.

    \subsubsection*{Groupes d'aliments}

    \begin{align*}
        \attr{alim\_grp\_code}     &\rightarrow \attr{alim\_grp\_nom\_fr},\; \attr{alim\_grp\_nom\_eng} \\
        \attr{alim\_ssgrp\_code}   &\rightarrow \attr{alim\_ssgrp\_nom\_fr},\; \attr{alim\_ssgrp\_nom\_eng} \\
        \attr{alim\_ssssgrp\_code} &\rightarrow \attr{alim\_ssssgrp\_nom\_fr},\; \attr{alim\_ssssgrp\_nom\_eng}
    \end{align*}

    \subsubsection*{Aliments}

    \begin{align*}
        \attr{alim\_code} &\rightarrow \attr{alim\_nom\_fr},\; \attr{alim\_nom\_eng},\; \attr{alim\_nom\_sci},\; \attr{facteur\_jones} \\
        \attr{alim\_code} &\rightarrow \attr{alim\_grp\_code},\; \attr{alim\_ssgrp\_code},\; \attr{alim\_ssssgrp\_code}
    \end{align*}

    \subsubsection*{Constituant}

    \[
        \attr{const\_code} \rightarrow \attr{const\_nom\_fr},\; \attr{const\_nom\_eng},\; \attr{code\_infoods}
    \]

    \subsubsection*{Composition}

    \[
        (\attr{alim\_code},\, \attr{const\_code}) \rightarrow \attr{teneur},\; \attr{min},\; \attr{max},\; \attr{code\_confiance},\; \attr{source\_code}
    \]

    \subsubsection*{Sources}

    \[
        \attr{source\_code} \rightarrow \attr{ref\_citation}
    \]

% ------------------------------------------------------------------

    \subsection{Dépendances fonctionnelles de \texorpdfstring{$V$}{V}}

    À partir de la documentation fournie et du fichier Excel, on déduit les dépendances fonctionnelles
    afin de constituer un schéma.

    On dispose des propriétés suivantes sur les ensembles de codes :

    \begin{itemize}
        \item $(\attr{alim\_moy\_code} \cup \attr{alim\_contrib\_code}) \subset \attr{alim\_code}$\\
        L'utilisation du symbole $\subseteq$ est exclue au stade des données brutes : on ne peut pas
        garantir que tout \attr{alim\_code} est référencé en tant qu'aliment moyen ou contributeur.
        Cette propriété ne sera satisfaite qu'après extension de la relation (voir
        section~\ref{sec:unifier}).

        \item $\attr{alim\_contrib\_code} \cap \attr{alim\_moy\_code} = \emptyset$
    \end{itemize}

    \subsubsection*{Aliments moyens}

    \begin{align*}
        \attr{alim\_moy\_code}    &\rightarrow \attr{alim\_moy\_nom} \\
        \attr{alim\_contrib\_code}&\rightarrow \attr{alim\_contrib\_nom} \\
        (\attr{alim\_moy\_code},\, \attr{alim\_contrib\_code}) &\rightarrow \attr{pourcentage}
    \end{align*}

% ------------------------------------------------------------------

    \subsection{Remarques}
    \label{sec:remarques}

    \subsubsection*{Première remarque --- Redondance entre \texorpdfstring{$U$}{U} et \texorpdfstring{$V$}{V}}
    \label{sec:remarque1}

    On constate une équivalence entre certaines dépendances fonctionnelles des aliments moyens et celles
    de la table CIQUAL. En effet :
    \[
        \attr{alim\_moy\_nom} \equiv \attr{alim\_contrib\_nom} \equiv \attr{alim\_nom\_fr}
    \]
    Les dépendances
    \begin{align*}
        \attr{alim\_moy\_code}    &\rightarrow \attr{alim\_moy\_nom} \\
        \attr{alim\_contrib\_code}&\rightarrow \attr{alim\_contrib\_nom}
    \end{align*}
    sont donc équivalentes à $\attr{alim\_code} \rightarrow \attr{alim\_nom\_fr}$ et produisent une
    redondance. Elles ont été inscrites précédemment afin de respecter la structure des fichiers sources
    originaux.

    \subsubsection*{Deuxième remarque --- Hiérarchie des groupes}
    \label{sec:remarque2}

    Après observation approfondie des données, on relève les dépendances additionnelles suivantes :
    \begin{align*}
        \attr{alim\_ssssgrp\_code} &\rightarrow \attr{alim\_ssgrp\_code} \\
        \attr{alim\_ssgrp\_code}   &\rightarrow \attr{alim\_grp\_code}
    \end{align*}

    \textbf{Note :} ces dépendances sont valables à condition de traiter tout groupe absent (groupe,
    sous-groupe ou sous-sous-groupe) comme \texttt{NULL}, et non comme porteur du code \texttt{000000}
    associé à un nom \texttt{NULL}.

    Deux solutions sont envisageables :

    \begin{enumerate}
        \item \textbf{Décomposition hiérarchique.} Diviser les trois types de groupe en tables distinctes,
        et faire référer chaque groupe fils à son groupe père par une cléf étrangère. Cette solution
        implique une micro-redondance sur les codes \texttt{NULL} (\texttt{000000}, \texttt{0000}),
        mais uniquement pour les codes, les noms étant \texttt{NULL}.\\[0.3em]
        \textbf{Cependant, cette solution viole la Deuxième Forme Normale (2NF).}
        En effet, si la cléf d'un sous-groupe est définie comme le couple
        $(\attr{code\_fils},\, \attr{code\_pere})$, alors $\attr{alim\_ssgroupe\_nom}$ dépend
        uniquement de $\attr{code\_fils}$ : il s'agit d'une \emph{dépendance partielle} au sens de
        la Notion~XVII.10 \hspace{0.2em}\citeg{Gardarin, p.\,690--691}, ce qui constitue une
        violation caractérisée de la 2NF.

        \item \textbf{Conservation de la structure XML.} Définir chaque enregistrement de groupe par le
        triplet $(\attr{alim\_grp\_code},\, \attr{alim\_ssgrp\_code},\, \attr{alim\_ssssgrp\_code})$,
        qui détermine les noms. Chaque groupe est par ailleurs défini indépendamment par son code,
        ce qui élimine la redondance des noms. La hiérarchie des groupes est portée par le triplet.
        Cette solution traite également naturellement le cas \texttt{NULL} mentionné ci-dessus.
    \end{enumerate}

    \textbf{La solution retenue est la deuxième}, car la première viole la 2NF comme expliqué ci-dessus.
    De plus, en conception de bases de données, les codes de sous-groupe et de sous-sous-groupe ne doivent
    pas être construits par agrégation (fusion) des codes des groupes parents.

% ==================================================================


    \section{Cléf de \texorpdfstring{$U$}{U}}

    \begin{rappelbox}
        \textbf{Notion XVII.8 : Cléf de relation (\textit{Relation key})}
        \hfill\citeg{Gardarin, p.\,689}\\[0.5em]
        Sous-ensemble $X$ des attributs d'une relation $R(A_1, A_2, \ldots, A_n)$ tel que :
        \begin{enumerate}
            \item $X \rightarrow A_1\, A_2\, \ldots\, A_n$
            \item Il n'existe pas de sous-ensemble $Y \subsetneq X$ tel que $Y \rightarrow A_1\, A_2\, \ldots\, A_n$.
        \end{enumerate}
        En clair, une cléf est un ensemble \emph{minimal} d'attributs qui détermine tous les autres.
        Un ensemble d'attributs qui inclut une cléf sans en être une est appelé \emph{supercléf}. On parle
        de \emph{cléf candidate} pour désigner une cléf quelconque, et de \emph{cléf primaire} pour la cléf
        choisie comme identifiant principal.
    \end{rappelbox}

    D'après les dépendances fonctionnelles établies, le couple $(\attr{alim\_code},\, \attr{const\_code})$
    est une cléf minimale de $U$ : il est irréductible et détermine l'ensemble de tous les attributs de $U$.

% ==================================================================


    \section{Cléf de \texorpdfstring{$V$}{V}}

    De même, le couple $(\attr{alim\_moy\_code},\, \attr{alim\_contrib\_code})$ est une cléf minimale de
    $V$. Ce couple constitue, par construction, un couple d'instances de \attr{alim\_code}, chacun
    désignant un aliment distinct : l'un générique (aliment moyen), l'autre contributeur. Il est
    irréductible et détermine l'ensemble de tous les attributs de $V$.

% ==================================================================


    \section{Unifier \texorpdfstring{$U$}{U} et \texorpdfstring{$V$}{V}}
    \label{sec:unifier}

    Pour définir un schéma unifié de $U$ et $V$, il faut satisfaire la contrainte :
    \[
        (\attr{alim\_moy\_code} \cup \attr{alim\_contrib\_code}) \subseteq \attr{alim\_code}
    \]

    Les données brutes vérifient seulement l'inclusion stricte $\subset$ (environ 40\,\% des
    \attr{alim\_code} sont référencés comme aliment moyen ou contributeur). Afin d'étendre cette relation
    à l'ensemble des \attr{alim\_code}, on procède à l'insertion de \emph{n-uplets artificiels} pour tout
    $c \in \attr{alim\_code}$ tel que
    $c \notin (\attr{alim\_moy\_code} \cup \attr{alim\_contrib\_code})$ :

    \begin{itemize}
        \item $\attr{alim\_moy\_code} = \texttt{NULL}$
        \item $\attr{alim\_contrib\_code} = c$ \quad (code présent dans \attr{alim\_code} mais non encore référencé)
        \item $\attr{pourcentage} = \texttt{NULL}$
    \end{itemize}

    \begin{notebox}
        \small
        Ces n-uplets artificiels ont pour seul objet de satisfaire formellement la contrainte d'inclusion
        ensembliste. Ils ne portent aucune information sémantique sur une appartenance réelle à un groupe
        d'aliments moyens ou de contributeurs. Leur présence dans la relation est une décision de conception
        et non un fait du domaine.
    \end{notebox}

    Après cette extension, on obtient :
    \[
        (\attr{alim\_moy\_code} \cup \attr{alim\_contrib\_code}) = \attr{alim\_code}
    \]
    La disjonction $\attr{alim\_contrib\_code} \cap \attr{alim\_moy\_code} = \emptyset$ reste satisfaite
    puisque les n-uplets artificiels sont insérés uniquement dans \attr{alim\_contrib\_code}.

    Ainsi, \attr{alim\_moy\_code} contient l'ensemble des aliments génériques, et
    \attr{alim\_contrib\_code} contient l'ensemble des aliments susceptibles d'être ou non associés à un
    aliment générique.

% ==================================================================


    \section{Algorithme de Bernstein}

    \begin{rappelbox}
        \textbf{Algorithme de synthèse en Troisième Forme Normale}
        \hfill\citeg{Gardarin, p.\,695}\\[0.5em]
        Une décomposition en 3\textsuperscript{e} FN peut être générée par un algorithme de synthèse ayant
        pour entrées l'ensemble des attributs ainsi que les dépendances fonctionnelles. Le principe consiste à
        construire d'abord une \emph{couverture minimale} $F$ des DF élémentaires, puis à partitionner $F$
        en groupes $F_i$ tels que les DF de chaque $F_i$ aient le même ensemble d'attributs à gauche. Chaque
        groupe produit une relation en 3\textsuperscript{e} FN dont la cléf est cet ensemble d'attributs
        gauches.
    \end{rappelbox}

    \begin{rappelbox}
        \textbf{Notion XVII.7 : Couverture minimale (\textit{Minimal cover})}
        \hfill\citeg{Gardarin, p.\,688}\\[0.5em]
        Ensemble $F$ de DF élémentaires associé à un ensemble d'attributs vérifiant les deux propriétés
        suivantes :
        \begin{enumerate}
            \item Aucune dépendance dans $F$ n'est redondante : pour toute DF $f \in F$, l'ensemble $F \setminus \{f\}$ n'est pas équivalent à $F$ (sa fermeture transitive est strictement plus petite).
            \item Toute DF élémentaire de l'ensemble des attributs appartient à la fermeture transitive $F^+$.
        \end{enumerate}
        La couverture minimale n'est en général pas unique.
    \end{rappelbox}

    \begin{rappelbox}
        \textbf{Référence algorithmique}\\[0.3em]
        \textit{Bernstein P.A., « Synthesizing Third Normal Form Relations from Functional Dependencies »,
            ACM Transactions on Database Systems, Vol.\,1, n°\,4, p.\,277--298, 1976.}
    \end{rappelbox}

% ------------------------------------------------------------------

    \subsection{Couverture minimale}

    \subsubsection*{Réduction à droite}
    \bigskip
    \begin{rappelbox}
        \textbf{Notion XVII.6 : Dépendance fonctionnelle élémentaire}
        \hfill\citeg{Gardarin, p.\,686}\\[0.5em]
        Dépendance fonctionnelle de la forme $X \rightarrow A$, où $A$ est un attribut unique n'appartenant
        pas à $X$ ($A \notin X$) et où il n'existe pas $X' \subsetneq X$ tel que $X' \rightarrow A$.\\[0.3em]
        La seule règle d'inférence s'appliquant aux DF élémentaires est la transitivité.
    \end{rappelbox}

    La réduction à droite consiste à décomposer chaque DF de façon à n'avoir qu'un seul attribut à
    droite, conformément à la Notion~XVII.6 ci-dessus.

    \begin{align*}
        \attr{alim\_grp\_code}     &\rightarrow \attr{alim\_grp\_nom\_fr} \\
        \attr{alim\_grp\_code}     &\rightarrow \attr{alim\_grp\_nom\_eng} \\
        \attr{alim\_ssgrp\_code}   &\rightarrow \attr{alim\_ssgrp\_nom\_fr} \\
        \attr{alim\_ssgrp\_code}   &\rightarrow \attr{alim\_ssgrp\_nom\_eng} \\
        \attr{alim\_ssssgrp\_code} &\rightarrow \attr{alim\_ssssgrp\_nom\_fr} \\
        \attr{alim\_ssssgrp\_code} &\rightarrow \attr{alim\_ssssgrp\_nom\_eng} \\
        \attr{alim\_code}          &\rightarrow \attr{alim\_nom\_fr} \\
        \attr{alim\_code}          &\rightarrow \attr{alim\_nom\_eng} \\
        \attr{alim\_code}          &\rightarrow \attr{alim\_nom\_sci} \\
        \attr{alim\_code}          &\rightarrow \attr{facteur\_jones} \\
        \attr{alim\_code}          &\rightarrow \attr{alim\_grp\_code} \\
        \attr{alim\_code}          &\rightarrow \attr{alim\_ssgrp\_code} \\
        \attr{alim\_code}          &\rightarrow \attr{alim\_ssssgrp\_code} \\
        \attr{const\_code}         &\rightarrow \attr{const\_nom\_fr} \\
        \attr{const\_code}         &\rightarrow \attr{const\_nom\_eng} \\
        \attr{const\_code}         &\rightarrow \attr{code\_infoods} \\
        (\attr{alim\_code},\,\attr{const\_code}) &\rightarrow \attr{teneur} \\
        (\attr{alim\_code},\,\attr{const\_code}) &\rightarrow \attr{min} \\
        (\attr{alim\_code},\,\attr{const\_code}) &\rightarrow \attr{max} \\
        (\attr{alim\_code},\,\attr{const\_code}) &\rightarrow \attr{code\_confiance} \\
        (\attr{alim\_code},\,\attr{const\_code}) &\rightarrow \attr{source\_code} \\
        \attr{source\_code}        &\rightarrow \attr{ref\_citation} \\
        \attr{alim\_moy\_code}     &\rightarrow \attr{alim\_moy\_nom} \\
        \attr{alim\_contrib\_code} &\rightarrow \attr{alim\_contrib\_nom} \\
        (\attr{alim\_moy\_code},\,\attr{alim\_contrib\_code}) &\rightarrow \attr{pourcentage}
    \end{align*}

    \subsubsection*{Réduction à gauche}

    Par application de la remarque~\ref{sec:remarque1} (équivalence sémantique des noms), les dépendances
    redondantes suivantes sont remplacées :
    \begin{align*}
        \attr{alim\_moy\_code}    &\rightarrow \attr{alim\_moy\_nom}
        \quad \longmapsto \quad \attr{alim\_code} \rightarrow \attr{alim\_nom\_fr} \\
        \attr{alim\_contrib\_code}&\rightarrow \attr{alim\_contrib\_nom}
        \quad \longmapsto \quad \attr{alim\_code} \rightarrow \attr{alim\_nom\_fr}
    \end{align*}

    \subsubsection*{Suppression des redondances}

    Après réduction à droite, réduction à gauche et élimination des DF redondantes, on obtient la
    couverture minimale suivante :

    \begin{align*}
        \attr{alim\_grp\_code}     &\rightarrow \attr{alim\_grp\_nom\_fr} \\
        \attr{alim\_grp\_code}     &\rightarrow \attr{alim\_grp\_nom\_eng} \\
        \attr{alim\_ssgrp\_code}   &\rightarrow \attr{alim\_ssgrp\_nom\_fr} \\
        \attr{alim\_ssgrp\_code}   &\rightarrow \attr{alim\_ssgrp\_nom\_eng} \\
        \attr{alim\_ssssgrp\_code} &\rightarrow \attr{alim\_ssssgrp\_nom\_fr} \\
        \attr{alim\_ssssgrp\_code} &\rightarrow \attr{alim\_ssssgrp\_nom\_eng} \\
        \attr{alim\_code}          &\rightarrow \attr{alim\_nom\_fr} \\
        \attr{alim\_code}          &\rightarrow \attr{alim\_nom\_eng} \\
        \attr{alim\_code}          &\rightarrow \attr{alim\_nom\_sci} \\
        \attr{alim\_code}          &\rightarrow \attr{facteur\_jones} \\
        \attr{alim\_code}          &\rightarrow \attr{alim\_grp\_code} \\
        \attr{alim\_code}          &\rightarrow \attr{alim\_ssgrp\_code} \\
        \attr{alim\_code}          &\rightarrow \attr{alim\_ssssgrp\_code} \\
        \attr{const\_code}         &\rightarrow \attr{const\_nom\_fr} \\
        \attr{const\_code}         &\rightarrow \attr{const\_nom\_eng} \\
        \attr{const\_code}         &\rightarrow \attr{code\_infoods} \\
        (\attr{alim\_code},\,\attr{const\_code}) &\rightarrow \attr{teneur} \\
        (\attr{alim\_code},\,\attr{const\_code}) &\rightarrow \attr{min} \\
        (\attr{alim\_code},\,\attr{const\_code}) &\rightarrow \attr{max} \\
        (\attr{alim\_code},\,\attr{const\_code}) &\rightarrow \attr{code\_confiance} \\
        (\attr{alim\_code},\,\attr{const\_code}) &\rightarrow \attr{source\_code} \\
        \attr{source\_code}        &\rightarrow \attr{ref\_citation} \\
        (\attr{alim\_moy\_code},\,\attr{alim\_contrib\_code}) &\rightarrow \attr{pourcentage}
    \end{align*}

% ------------------------------------------------------------------

    \subsection{Regroupement des DF et création des schémas}

    \begin{rappelbox}
        \textbf{Notion XVII.11 : Troisième forme normale (\textit{Third normal form})}
        \hfill\citeg{Gardarin, p.\,691--692}\\[0.5em]
        Une relation $R$ est en troisième forme normale si et seulement si :
        \begin{enumerate}
            \item Elle est en deuxième forme normale.
            \item Tout attribut n'appartenant pas à une cléf ne dépend pas d'un autre attribut non-cléf.
        \end{enumerate}
    \end{rappelbox}

    \begin{notebox}
        \small
        \textbf{Précision.} La 3\textsuperscript{e} FN implique automatiquement la 2\textsuperscript{e} FN :
        toute relation en 3NF est nécessairement en 2NF (la condition 1 est automatiquement vérifiée dès
        lors que la condition 2 est satisfaite, mais elle figure par souci de clarté). Si la relation possède
        plusieurs cléfs candidates, la définition doit être vérifiée pour chacune d'elles
        successivement. \hfill\citeg{Gardarin, p.\,692}
    \end{notebox}

    On regroupe les dépendances fonctionnelles par déterminant pour constituer les relations en 3NF.
    Les attributs soulignés constituent la cléf primaire de chaque relation ; les contraintes d'inclusion
    ($\subseteq$) expriment les cléfs étrangères.

    \bigskip
    \noindent\textbf{Relations de groupes :}
    \begin{align*}
        &\textbf{groupe}(\underline{\attr{alim\_grp\_code}},\; \attr{alim\_grp\_nom\_fr},\; \attr{alim\_grp\_nom\_eng}) \\
        &\textbf{ssgroupe}(\underline{\attr{alim\_ssgrp\_code}},\; \attr{alim\_ssgrp\_nom\_fr},\; \attr{alim\_ssgrp\_nom\_eng}) \\
        &\textbf{ssssgroupe}(\underline{\attr{alim\_ssssgrp\_code}},\; \attr{alim\_ssssgrp\_nom\_fr},\; \attr{alim\_ssssgrp\_nom\_eng})
    \end{align*}

    Afin de satisfaire la remarque~\ref{sec:remarque2}, on définit la relation de hiérarchie suivante,
    qui porte la structure des groupes sous la forme d'un triplet-cléf :
    \[
        \textbf{hierarchie\_groupe}(\underline{\attr{alim\_grp\_code},\; \attr{alim\_ssgrp\_code},\; \attr{alim\_ssssgrp\_code}})
    \]

    avec les contraintes d'inclusion :
    \begin{align*}
        \attr{alim\_grp\_code}     &\subseteq \textbf{groupe}[\attr{alim\_grp\_code}] \\
        \attr{alim\_ssgrp\_code}   &\subseteq \textbf{ssgroupe}[\attr{alim\_ssgrp\_code}] \\
        \attr{alim\_ssssgrp\_code} &\subseteq \textbf{ssssgroupe}[\attr{alim\_ssssgrp\_code}]
    \end{align*}

    \bigskip
    \noindent\textbf{Relation des aliments :}
    \[
        \textbf{aliments}(\underline{\attr{alim\_code}},\; \attr{alim\_nom\_fr},\; \attr{alim\_nom\_eng},\; \attr{alim\_nom\_sci},\; \attr{facteur\_jones},\;
    \]
    \[ \attr{alim\_grp\_code},\; \attr{alim\_ssgrp\_code},\; \attr{alim\_ssssgrp\_code})\]
    \[
        (\attr{alim\_grp\_code},\; \attr{alim\_ssgrp\_code},\; \attr{alim\_ssssgrp\_code}) \subseteq \textbf{hierarchie\_groupe}[
    \]
    \[\attr{alim\_grp\_code},\; \attr{alim\_ssgrp\_code},\; \attr{alim\_ssssgrp\_code}]\]

    \bigskip
    \noindent\textbf{Relations auxiliaires :}
    \begin{align*}
        &\textbf{constituant}(\underline{\attr{const\_code}},\; \attr{const\_nom\_fr},\; \attr{const\_nom\_eng},\; \attr{code\_infoods}) \\
        &\textbf{sources}(\underline{\attr{source\_code}},\; \attr{ref\_citation})
    \end{align*}

    \bigskip
    \noindent\textbf{Relation de composition :}
    \[
        \textbf{composition}(\underline{\attr{alim\_code},\; \attr{const\_code}},\; \attr{teneur},\; \attr{min},\; \attr{max},\; \attr{code\_confiance},\; \attr{source\_code})
    \]
    \begin{align*}
        \attr{alim\_code}   &\subseteq \textbf{aliments}[\attr{alim\_code}] \\
        \attr{const\_code}  &\subseteq \textbf{constituant}[\attr{const\_code}] \\
        \attr{source\_code} &\subseteq \textbf{sources}[\attr{source\_code}]
    \end{align*}

    \bigskip
    \noindent\textbf{Relation des aliments moyens :}
    \[
        \textbf{alim\_moyen}(\underline{\attr{alim\_moy\_code},\; \attr{alim\_contrib\_code}},\; \attr{pourcentage})
    \]
    \begin{align*}
        \attr{alim\_moy\_code}    &\subseteq \textbf{aliments}[\attr{alim\_code}] \\
        \attr{alim\_contrib\_code}&\subseteq \textbf{aliments}[\attr{alim\_code}]
    \end{align*}

% ==================================================================


    \section{Création de la base de données}

    La base de données a été établie à partir de la conception effectuée précédemment.

% ------------------------------------------------------------------

    \subsection{Raffinement de l'attribut \texttt{teneur}}

    Pour l'attribut \attr{teneur} de la relation \textbf{composition}, on distingue trois types de valeurs
    présentes dans les fichiers sources :

    \begin{itemize}
        \item \textbf{Valeur exacte :} \texttt{14,9} — \texttt{1410} — \texttt{0,547}
        \item \textbf{Valeur inférieure à un seuil :} \texttt{< 700} — \texttt{< 0,001} — \texttt{< 256}
        \item \textbf{Valeur trace :} \texttt{traces}
    \end{itemize}
    \bigskip
    \begin{rappelbox}
        \textbf{Notion XVII.9 : Première forme normale (\textit{First normal form})}
        \hfill\citeg{Gardarin, p.\,689--690}\\[0.5em]
        Une relation est en première forme normale si tout attribut contient une \emph{valeur atomique}.
    \end{rappelbox}

    \begin{notebox}
        \small
        Le raffinement de \attr{teneur} relève d'un double impératif.\\[0.3em]
        \textbf{Impératif de modélisation (1NF).} L'attribut \attr{teneur\_brute} encode à la fois un
        qualificatif de type et une valeur numérique, voire un symbole non numérique (\texttt{traces}). Sa
        décomposition en attributs distincts rend le schéma conforme au principe d'atomicité de la
        1NF au sens applicatif.\\[0.3em]
        \textbf{Impératif de traitement analytique.} Le mot \texttt{traces} est une valeur non numérique
        qui bloque tout calcul statistique (moyenne, variance, etc.). Sa transformation en valeur
        conventionnelle $0$ dans \attr{teneur\_valeur} permet ces traitements, tandis que
        \attr{teneur\_brute} conserve l'information sémantique pour la présentation à l'utilisateur, où
        afficher $0$ à la place de \texttt{traces} serait sémantiquement incorrect.
    \end{notebox}

    En conséquence, \attr{teneur} est décomposé en trois attributs :

    \begin{itemize}
        \item \attr{teneur\_brute} : chaîne de caractères reproduisant la valeur telle qu'inscrite dans le
        fichier source.
        \item \attr{teneur\_type} : type de teneur du couple-cléf, prenant l'une des valeurs :
        \texttt{EXACTE}, \texttt{INFERIEURE}, \texttt{TRACE}.
        \item \attr{teneur\_valeur} : valeur numérique de la teneur, définie selon \attr{teneur\_type} :
        \begin{itemize}
            \item Pour \texttt{EXACTE} : la valeur numérique intacte.
            \item Pour \texttt{INFERIEURE} : la valeur numérique privée du symbole \texttt{<}.
            \item Pour \texttt{TRACE} : la valeur conventionnelle $0$.
        \end{itemize}
    \end{itemize}

% ==================================================================
    \newpage


    \section*{Références bibliographiques}
    \addcontentsline{toc}{section}{Références bibliographiques}

    Les notions théoriques mobilisées dans ce document s'appuient sur l'ouvrage suivant :

    \bigskip
    \noindent\textbf{[Gardarin]} \quad Gardarin, G. \textit{Bases de données}. Eyrolles, 1\textsuperscript{re} édition. 826~pages.

    \bigskip
    \noindent Les pages précisément citées sont les suivantes :

    \begin{itemize}
        \item \textbf{p.\,685} --- Notion~XVII.5 : Dépendance fonctionnelle. Définition formelle.
        Remarque sur le caractère intentionnel (et non extensionnel) des DF : impossibilité de les
        déduire d'une réalisation particulière de la relation.

        \item \textbf{p.\,686} --- Notion~XVII.6 : Dépendance fonctionnelle élémentaire. DF à partie
        droite singleton, irréductible à gauche. Axiomes d'Armstrong (réflexivité, augmentation,
        transitivité) et règles dérivées (union, pseudo-transitivité, décomposition).

        \item \textbf{p.\,688} --- Notion~XVII.7 : Couverture minimale. Définition par les propriétés
        d'absence de redondance et de complétude par fermeture transitive. Non-unicité de la
        couverture minimale.

        \item \textbf{p.\,689} --- Notion~XVII.8 : cléf de relation. Définition formelle par les DF.
        Distinction cléf primaire / supercléf / cléf candidate.

        \item \textbf{p.\,689--690} --- Notion~XVII.9 : Première forme normale. Atomicité des valeurs
        d'attributs et opération de désimbrication (UNNEST).

        \item \textbf{p.\,690--691} --- Notion~XVII.10 : Deuxième forme normale. Absence de dépendance
        partielle d'un attribut non-cléf sur une partie de la cléf. Schéma typique de violation et
        règle de décomposition associée.

        \item \textbf{p.\,691--692} --- Notion~XVII.11 : Troisième forme normale. Absence de dépendance
        transitive entre attributs non-cléf. Implication de la 3NF par la 2NF. Vérification pour
        chaque cléf candidate.

        \item \textbf{p.\,695} --- Algorithme de synthèse en 3\textsuperscript{e} FN. Principe de
        construction par couverture minimale et partitionnement en groupes de DF de même déterminant.
        Chaque groupe produit une relation en 3NF dont la cléf est le déterminant commun.
    \end{itemize}

    \bigskip
    \noindent\textbf{[Bernstein76]} \quad Bernstein, P.A. \textit{« Synthesizing Third Normal Form
    Relations from Functional Dependencies »}. ACM Transactions on Database Systems, Vol.\,1, n°\,4,
    p.\,277--298, 1976.\\
    \textit{Algorithme de synthèse appliqué en section~\ref{sec:unifier} et suivantes pour la
    génération du schéma en 3NF à partir de la couverture minimale.}

    \bigskip
    \noindent\textbf{[Ciqual2025]} \quad Du~Chaffaut,~L., Oseredczuk,~M., Gauvreau-Béziat,~J.
    \textit{Table de composition nutritionnelle des aliments Ciqual 2025}.
    Version~1. Agence nationale de sécurité sanitaire de l'alimentation, de l'environnement
    et du travail (ANSES). Recherche Data Gouv, 2025.\\
    \textsc{doi} : \href{https://doi.org/10.57745/RDMHWY}{\texttt{10.57745/RDMHWY}}
    [consulté le 19~novembre~2025].

\end{document}